\documentclass{amsart}
%	include packages
\usepackage{amssymb}
\usepackage{graphicx}
\usepackage[cmtip,all]{xy}
\usepackage{bbm}
\usepackage{comment}
\usepackage{mathrsfs}
%\usepackage{eucal}
\usepackage{microtype}
\usepackage{xcolor}
\usepackage{amscd,verbatim}
\usepackage{hyperref}
\hypersetup{
 colorlinks,
 linkcolor={blue!60!black},
 citecolor={blue!60!black},
 urlcolor={blue!60!black}
}





%	page layout
\oddsidemargin =10mm
\evensidemargin =10mm
\topmargin =5mm
\textwidth =150mm
\textheight =220mm
\calclayout








%	define "define symbol/category" commands
\def\DefineSymbol#1#2{\newcommand{#1}{{\mathrm {#2}}}}
\def\DefineCategory#1#2{\newcommand{#1}{{\mathrm {#2}}}}





%	define symbols
\DefineSymbol{\pr}{pr}
\DefineSymbol{\id}{id}
\DefineSymbol{\const}{const}
\DefineSymbol{\op}{op}
\DefineSymbol{\diag}{diag}



%	define operators
\DeclareMathOperator{\Hom}{Hom}
\DeclareMathOperator{\Fun}{Fun}
\DeclareMathOperator{\Ker}{Ker}
\DeclareMathOperator{\Coker}{Coker}
\let\Im\relax
\DeclareMathOperator{\Im}{Im}
\DeclareMathOperator{\Aut}{Aut}
\DeclareMathOperator{\Gal}{Gal}
\DeclareMathOperator{\Frac}{Frac}
\DeclareMathOperator{\dlog}{dlog}
\DeclareMathOperator{\height}{ht}
\DeclareMathOperator*{\colim}{colim}
\DeclareMathOperator{\eq}{eq}
\DeclareMathOperator{\coeq}{coeq}
\DeclareMathOperator{\Ass}{Ass}
\DeclareMathOperator{\Nm}{Nm}
\DeclareMathOperator{\Bl}{Bl}
\DeclareMathOperator{\Pro}{Pro}
\DeclareMathOperator{\PSh}{PSh}
\DeclareMathOperator{\Shv}{Shv}
\DeclareMathOperator{\Sym}{Sym}
\DeclareMathOperator{\Spec}{Spec}
\DeclareMathOperator{\End}{End}



%	define categories
\DefineCategory{\Set}{Set}
\DefineCategory{\Ab}{Ab}
\DefineCategory{\Mod}{Mod}
\DefineCategory{\Alg}{Alg}
\DefineCategory{\Ch}{Ch}
\DefineCategory{\Mon}{Mon}
\DefineCategory{\CMon}{CMon}



%\mathbb
\newcommand{\Abb}{\mathbb{A}}
\renewcommand{\Bbb}{\mathbb{B}}
\newcommand{\Cbb}{\mathbb{C}}
\newcommand{\Dbb}{\mathbb{D}}
\newcommand{\Ebb}{\mathbb{E}}
\newcommand{\Fbb}{\mathbb{F}}
\newcommand{\Gbb}{\mathbb{G}}
\newcommand{\Hbb}{\mathbb{H}}
\newcommand{\Ibb}{\mathbb{I}}
\newcommand{\Jbb}{\mathbb{J}}
\newcommand{\Kbb}{\mathbb{K}}
\newcommand{\Lbb}{\mathbb{L}}
\newcommand{\Mbb}{\mathbb{M}}
\newcommand{\Nbb}{\mathbb{N}}
\newcommand{\Obb}{\mathbb{O}}
\newcommand{\Pbb}{\mathbb{P}}
\newcommand{\Qbb}{\mathbb{Q}}
\newcommand{\Rbb}{\mathbb{R}}
\newcommand{\Sbb}{\mathbb{S}}
\newcommand{\Tbb}{\mathbb{T}}
\newcommand{\Ubb}{\mathbb{U}}
\newcommand{\Vbb}{\mathbb{V}}
\newcommand{\Wbb}{\mathbb{W}}
\newcommand{\Xbb}{\mathbb{X}}
\newcommand{\Ybb}{\mathbb{Y}}
\newcommand{\Zbb}{\mathbb{Z}}



%\mathcal
\newcommand{\Acal}{\mathcal{A}}
\newcommand{\Bcal}{\mathcal{B}}
\newcommand{\Ccal}{\mathcal{C}}
\newcommand{\Dcal}{\mathcal{D}}
\newcommand{\Ecal}{\mathcal{E}}
\newcommand{\Fcal}{\mathcal{F}}
\newcommand{\Gcal}{\mathcal{G}}
\newcommand{\Hcal}{\mathcal{H}}
\newcommand{\Ical}{\mathcal{I}}
\newcommand{\Jcal}{\mathcal{J}}
\newcommand{\Kcal}{\mathcal{K}}
\newcommand{\Lcal}{\mathcal{L}}
\newcommand{\Mcal}{\mathcal{M}}
\newcommand{\Ncal}{\mathcal{N}}
\newcommand{\Ocal}{\mathcal{O}}
\newcommand{\Pcal}{\mathcal{P}}
\newcommand{\Qcal}{\mathcal{Q}}
\newcommand{\Rcal}{\mathcal{R}}
\newcommand{\Scal}{\mathcal{S}}
\newcommand{\Tcal}{\mathcal{T}}
\newcommand{\Ucal}{\mathcal{U}}
\newcommand{\Vcal}{\mathcal{V}}
\newcommand{\Wcal}{\mathcal{W}}
\newcommand{\Xcal}{\mathcal{X}}
\newcommand{\Ycal}{\mathcal{Y}}
\newcommand{\Zcal}{\mathcal{Z}}


%\mathfrak
\newcommand{\Afrak}{\mathfrak{A}}
\newcommand{\Bfrak}{\mathfrak{B}}
\newcommand{\Cfrak}{\mathfrak{C}}
\newcommand{\Dfrak}{\mathfrak{D}}
\newcommand{\Efrak}{\mathfrak{E}}
\newcommand{\Ffrak}{\mathfrak{F}}
\newcommand{\Gfrak}{\mathfrak{G}}
\newcommand{\Hfrak}{\mathfrak{H}}
\newcommand{\Ifrak}{\mathfrak{I}}
\newcommand{\Jfrak}{\mathfrak{J}}
\newcommand{\Kfrak}{\mathfrak{K}}
\newcommand{\Lfrak}{\mathfrak{L}}
\newcommand{\Mfrak}{\mathfrak{M}}
\newcommand{\Nfrak}{\mathfrak{N}}
\newcommand{\Ofrak}{\mathfrak{O}}
\newcommand{\Pfrak}{\mathfrak{P}}
\newcommand{\Qfrak}{\mathfrak{Q}}
\newcommand{\Rfrak}{\mathfrak{R}}
\newcommand{\Sfrak}{\mathfrak{S}}
\newcommand{\Tfrak}{\mathfrak{T}}
\newcommand{\Ufrak}{\mathfrak{U}}
\newcommand{\Vfrak}{\mathfrak{V}}
\newcommand{\Wfrak}{\mathfrak{W}}
\newcommand{\Xfrak}{\mathfrak{X}}
\newcommand{\Yfrak}{\mathfrak{Y}}
\newcommand{\Zfrak}{\mathfrak{Z}}


%\mathscr
\newcommand{\Ascr}{\mathscr{A}}
\newcommand{\Bscr}{\mathscr{B}}
\newcommand{\Cscr}{\mathscr{C}}
\newcommand{\Dscr}{\mathscr{D}}
\newcommand{\Escr}{\mathscr{E}}
\newcommand{\Fscr}{\mathscr{F}}
\newcommand{\Gscr}{\mathscr{G}}
\newcommand{\Hscr}{\mathscr{H}}
\newcommand{\Iscr}{\mathscr{I}}
\newcommand{\Jscr}{\mathscr{J}}
\newcommand{\Kscr}{\mathscr{K}}
\newcommand{\Lscr}{\mathscr{L}}
\newcommand{\Mscr}{\mathscr{M}}
\newcommand{\Nscr}{\mathscr{N}}
\newcommand{\Oscr}{\mathscr{O}}
\newcommand{\Pscr}{\mathscr{P}}
\newcommand{\Qscr}{\mathscr{Q}}
\newcommand{\Rscr}{\mathscr{R}}
\newcommand{\Sscr}{\mathscr{S}}
\newcommand{\Tscr}{\mathscr{T}}
\newcommand{\Uscr}{\mathscr{U}}
\newcommand{\Vscr}{\mathscr{V}}
\newcommand{\Wscr}{\mathscr{W}}
\newcommand{\Xscr}{\mathscr{X}}
\newcommand{\Yscr}{\mathscr{Y}}
\newcommand{\Zscr}{\mathscr{Z}}



%	use big tilde/hat/overline
\renewcommand{\tilde}{\widetilde}
\renewcommand{\hat}{\widehat}
\renewcommand{\bar}{\overline}



%	enumerate label (i),(ii),...
%\renewcommand{\labelenumi}{(\roman{enumi})}














\usepackage[style=alphabetic]{biblatex} 
\addbibresource{bibliography.bib}
\renewcommand\bibfont{\footnotesize}
\usepackage{scalerel}
%\usepackage{eucal}
\usepackage[nameinlink]{cleveref}

\oddsidemargin =10mm
\evensidemargin =10mm
\topmargin =5mm
\textwidth =150mm
\textheight =200mm
\calclayout


%	theorem environment
\theoremstyle{plain}
	\newtheorem{theorem}{Theorem}[section]
	\newtheorem{lemma}[theorem]{Lemma}
	\newtheorem{proposition}[theorem]{Proposition}
	\newtheorem{corollary}[theorem]{Corollary}
	\newtheorem{conjecture}[theorem]{Conjecture}
\theoremstyle{definition}
	\newtheorem{definition}[theorem]{Definition}
	\newtheorem{example}[theorem]{Example}
	\newtheorem{problem}[theorem]{Problem}
	\newtheorem{question}[theorem]{Question}
	\newtheorem{exercise}[theorem]{Exercise}
\theoremstyle{remark}
	\newtheorem{remark}[theorem]{Remark}
	\numberwithin{equation}{section}
	

\DeclareFieldFormat
  [article,book,inbook,incollection,inproceedings,patent,thesis,unpublished]
  {title}{\emph{#1\isdot}}
  

\newcommand{\bcube}{{\overline{\square}}}




\DeclareMathOperator{\Sh}{Sh}
\DeclareMathOperator{\Fil}{Fil}
\DeclareMathOperator{\Pic}{Pic}
\DeclareMathOperator{\trdeg}{trdeg}
\DeclareMathOperator{\Tr}{Tr}
\DeclareMathOperator{\RHom}{RHom}
\DeclareMathOperator{\codim}{codim}
\DeclareMathOperator{\cofib}{cofib}
\DeclareMathOperator{\Map}{Map}
\DeclareMathOperator{\Ext}{Ext}
\DeclareMathOperator{\CDiv}{CDiv}


\def\sm{\mathrm{sm}}
\def\tr{\mathrm{tr}}
\def\mtr{\mathrm{mtr}}
\def\smprop{\mathrm{sm,prop}}
\def\intHom{\mathcal{H}\mathrm{om}}
\def\ad{\mathrm{ad}}
\def\Suslin{\mathrm{S}}
\def\ceil#1{\lceil #1 \rceil}
\def\et{\mathrm{\acute{e}t}}
\def\Zar{\mathrm{Zar}}
\def\Nis{\mathrm{Nis}}
\def\eff{\mathrm{eff}}
\def\div{\mathrm{div}}
\def\D{\mathcal{D}}
\def\TBA{{\color{red}TBA}}
\def\LINE{{\color{red}\hrulefill}}
\def\prop{\mathrm{prop}}
\def\fin{\mathrm{fin}}
\def\ch{\mathrm{ch}}
\def\ls{\mathrm{ls}}
\def\HA{\mathrm{HA}}
\def\mot{\mathrm{mot}}
\def\red{\mathrm{red}}
\def\MOmega{\underline{\mathrm{M}}\Omega}
\def\MW{\underline{\mathrm{M}}\mathrm{W}}
\makeatother

\def\smallcube{{\scaleobj{0.8}{\overline{\square}}}}
\def\Lcube{h_0^{\smallcube}}
\def\Rcube{h^{0,\smallcube}}
\def\Lh0{h_0}
\def\Rh0{h^0}
\def\logrm{\mathrm{log}}
\def\nonlog{\mathrm{nonlog}}
\def\mod{\mathrm{mod}}
\def\bl{\mathrm{bl}}
\def\sNis{\mathrm{sNis}}
\def\dNis{\mathrm{dNis}}
\def\MM{\mathrm{MM}}

\DefineCategory{\Vect}{Vect}
\DefineCategory{\qCor}{qCor}
\DefineCategory{\qSm}{qSm}
\DefineCategory{\QCor}{QCor}
\DefineCategory{\QSm}{QSm}
\DefineCategory{\mCor}{mCor}
\DefineCategory{\mSm}{mSm}
\DefineCategory{\MCor}{MCor}
\DefineCategory{\MSm}{MSm}
\DefineCategory{\Ring}{Ring}
\DefineCategory{\Mor}{Mor}
\DefineCategory{\Sch}{Sch}
\DefineCategory{\Sm}{Sm}
\DefineCategory{\Cor}{Cor}
\DefineCategory{\DM}{DM}
\DefineCategory{\MDM}{MDM}
\DefineCategory{\mDM}{mDM}
\DefineCategory{\ulMSm}{\underline{M}Sm}
\DefineCategory{\ulMCor}{\underline{M}Cor}
\DefineCategory{\ulMDM}{\underline{M}DM}
\DefineCategory{\ulMNST}{\underline{M}NST}

\makeatletter
\def\quotprojlim{%
  \mathop{``\mathpalette\varlim@{\leftarrowfill@\textstyle}\hbox to 0pt{$"$}}\nmlimits@
}
\def\quotinjlim{%
  \mathop{``\mathpalette\varlim@{\rightarrowfill@\textstyle}"}\nmlimits@
}
\makeatother



\begin{document}
\title{Representability of cohomology theories with ramification filtration}

\author[J. Koizumi]{Junnosuke Koizumi}
\address{Graduate School of Mathematical Sciences, University of Tokyo, 3-8-1 Komaba, Meguro-ku, Tokyo 153-8914, Japan}
\email{jkoizumi@ms.u-tokyo.ac.jp}


\date{\today}
\thanks{}


\begin{abstract}
	In this paper, we provide a general criterion for a cohomology theory to be representable in the category of motives with modulus.
	Using this result, we prove the representability of the Witt vector cohomology under resolution of singularities.
\end{abstract}

%\subjclass[2010]{19E15 (14F42, 19D45, 19F15)}

\maketitle
\setcounter{tocdepth}{1}
\tableofcontents

\enlargethispage*{20pt}


\section{Introduction}

The category of motives with modulus $\ulMDM^\eff_k$ over a perfect field $k$ was defined by Kahn-Miyazaki-Saito-Yamazaki \cite{KMSY3} as an extension of Voevodsky's category of motives $\DM^\eff_k$.
It is expected to have the ability to capture many non-$\mathbb{A}^1$-invariant cohomology theories which do not factor through $\DM^\eff_k$.
Recently, Kelly-Miyazaki \cite{KellyMiyazaki_Hodge1} \cite{KellyMiyazaki_Hodge2} studied the Hodge cohomology $X\mapsto \mathrm{R}\Gamma(X,\Omega^q_X)$ with a ramification filtration and showed that it is representable in $\ulMDM^\eff_k$ when $\ch(k)=0$.
In this paper, we provide a general criterion for a cohomology theory to be representable in $\ulMDM^\eff_k$.
Using this result, we give an alternative proof of the representability of the Hodge cohomology due to Kelly-Miyazaki, and show that the Witt vector cohomology $X\mapsto \mathrm{R}\Gamma(X,\mathrm{W}_n\mathcal{O}_X)$ is also representable in $\ulMDM^\eff_k$ under resolution of singularities.

Just as $\DM^\eff_k$ is constructed from the category of smooth $k$-schemes, $\ulMDM^\eff_k$ is constructed from the category of \emph{modulus pairs}.
A modulus pair over $k$ is a pair $\mathcal{X}=(X,D_X)$ where $X$ is an algebraic $k$-scheme and $D_X$ is an effective divisor on $X$ such that $X\setminus|D_X|$ is smooth over $k$.
It is called \emph{log-smooth} if $X$ is smooth over $k$ and $|D_X|$ has strict normal crossings.
There is a category $\ulMCor_k$ of modulus pairs which is similar to Suslin-Voevodsky's category of finite correspondences $\Cor_k$.
We write $\ulMCor_k^\ls$ for the full subcategory of $\ulMCor_k$ spanned by log-smooth modulus pairs.
The modulus pair $\bcube=(\mathbb{P}^1,[\infty])$ is called the \emph{cube} and plays the role of $\mathbb{A}^1$ in Voevodsky's theory.

We say that a presheaf $F\colon (\ulMCor_k)^\op\to \Ab$ is a \emph{modulus (Nisnevich) sheaf} if for any $\mathcal{X}\in \ulMCor_k$, the presheaf $F_\mathcal{X}\colon U\mapsto F(U,D_X|_U)$ on $X_\Nis$ is a sheaf.
We write $\Sh_\Nis(\ulMCor_k)$ for the category of modulus sheaves.\footnote{This category is denoted by $\ulMNST_k$ in \cite{KMSY3}.}
For example, Kelly-Miyazaki \cite{KellyMiyazaki_Hodge2} constructed a modulus sheaf $\MOmega^q$ for $q\geq 0$ which satisfies
$$
	\MOmega^q(\mathcal{X})=\Gamma\bigl(X,\Omega^q_X(\log|D_X|)(D_X-|D_X|)\bigr)
$$
for $\mathcal{X}\in \ulMCor_k^\ls$.
If $\ch(k)=p>0$, we can construct a modulus sheaf $\MW_n$ for $n>0$ which satisfies
$$
	\MW_n(\mathcal{X})=\Gamma\biggl(X,\mathrm{W}_n\mathcal{O}_X\biggl(\frac{D_X-|D_X|}{p^{n-1}}\biggr)\biggr)
$$
for $\mathcal{X}\in \ulMCor_k^\ls$.
Here, $\mathrm{W}_n\mathcal{O}_X(D)$ denotes the \emph{Witt divisorial sheaf} introduced by Tanaka \cite{Tanaka_Witt}.
We consider the following conditions for a modulus sheaf $F$.
\begin{enumerate}
	\item	(semipurity, \cite[Definition 1.28]{Sai20}) $F(\mathcal{X})\to F(X^\circ)$ is injective for any $\mathcal{X}\in \ulMCor_k$.
	\item	(M-reciprocity, \cite[Definition 1.26]{Sai20}) $F$ is left Kan extended from the full subcategory of proper modulus pairs.
	\item	(ls cube invariance, \cite[Definition 1.12]{Sai20}) For any $\mathcal{X}\in \ulMCor_k^\ls$, the map $\pr_1^*\colon F(\mathcal{X})\to F(\mathcal{X}\otimes\bcube)$ is an isomorphism.
	\item	(cohomological ls cube invariance) For any $\mathcal{X}\in \ulMCor_k^\ls$, the morphism $F_\mathcal{X}\to \mathrm{R}\pr_{1,*}F_{\mathcal{X}\otimes \bcube}$ is an equivalence.
	\item	(affine vanishing property) For any effective divisors $E_1,\dots,E_n$ on $\mathbb{A}^1$ and $\mathcal{X}\in \ulMCor_k^\ls$, we have $\mathrm{R}^i\pr_{1,*}F_{\mathcal{X}\otimes(\mathbb{A}^1,E_1)\otimes\cdots\otimes(\mathbb{A}^1,E_n)}=0$ for $i>0$.
	\item	(cohomological blow-up invariance) For any $\mathcal{X}\in \ulMCor_k^\ls$ and any SNC blow-up (\Cref{def:SNC_blowup}) $\varphi\colon \mathcal{Y}\to \mathcal{X}$, the morphism $F_\mathcal{X}\to \mathrm{R}\varphi_* F_{\mathcal{Y}}$ is an equivalence.
\end{enumerate}
If (4) and (6) are satisfied, then (assuming resolution of singularities) we get an equivalence
$$
	\Map_{\ulMDM^\eff_k}\bigl(M(\mathcal{X}),F\bigr) \simeq \mathrm{R}\Gamma(X,F_\mathcal{X}).
$$
for any $\mathcal{X}\in \ulMCor_k^\ls$ (see \Cref{general_realization}).
In other words, (4) and (6) ensures that the cohomology of $F$ is representable in $\ulMDM^\eff_k$.
Saito proved (1)+(2)+(3)$\implies$(4) \cite[Theorem 0.6]{Sai20}, and Binda-R\"ulling-Saito proved (1)+(2)+(3)$\implies$(5) \cite[Lemma 2.10]{BRS}.
On the other hand, there is no known criterion for (6).
R\"ulling-Saito gave an example of a modulus sheaf satisfying (1)-(5) which does not satisfy (6) \cite[{}6.9]{RSram3}.

In order to find a good sufficient condition for (6), we enlarge $\ulMCor_k$ to a category $\QCor_k$ by allowing $\mathbb{Q}$-divisors in the definition of modulus pairs.
This idea was first introduced by the author and Miyazaki \cite{Koizumi-Miyazaki} to give a motivic construction of the de Rham-Witt complex.
We can consider conditions (1)-(6) also for modulus sheaves on $\QCor_k$.
After this enlargement, we can consider the following condition:
\begin{enumerate}\setcounter{enumi}{6}
	\item	(left continuity) The canonical map
	$\colim_{\varepsilon>0} F\bigl(X,(1-\varepsilon)D_X\bigr)\to F(\mathcal{X})$
	is an isomorphism for any $\mathcal{X}\in \QCor_k^\ls$.
\end{enumerate}
For example, the natural extensions of $\MOmega^q$ and $\MW_n$ to $\QCor_k$ have left continuity.
In \Cref{section_main}, we prove the following general criterion.

\begin{theorem}\label{main}
	Suppose that $F\in \Sh_\Nis(\QCor_k)$ has cohomological ls cube invariance, affine vanishing property, and left continuity.
	Then $F$ has cohomological blow-up invariance.
	In other words, we have (4)+(5)+(7)$\implies$(6).
\end{theorem}

Applying this to $\MOmega^q$ and $\MW_n$, we get the following non-trivial result.

\begin{corollary}\label{omega_witt_CBI}
	Let $\mathcal{X}\in \ulMCor_k^\ls$ and let $\varphi\colon \mathcal{Y}\to \mathcal{X}$ be an SNC blow-up.
	\begin{enumerate}
		\item	For any $i\geq 0$, $\varphi$ induces an isomorphism
				$\mathrm{H}^i\bigl(X,\MOmega^q_\mathcal{X})\xrightarrow{\sim} \mathrm{H}^i\bigl(Y,\MOmega^q_\mathcal{Y}\bigr)$.
		\item	If $\ch(k)=p>0$, then for any $i\geq 0$, $\varphi$ induces an isomorphism
				$\mathrm{H}^i\bigl(X,(\MW_n)_\mathcal{X}\bigr)
					\xrightarrow{\sim}
					\mathrm{H}^i\bigl(Y,(\MW_n)_\mathcal{Y}\bigr)$.
	\end{enumerate}
\end{corollary}



In order to represent these cohomology theories in $\ulMDM^\eff_k$, we assume resolution of singularities.

\begin{definition}\label{def:resolution}
	We say that $k$ admits \emph{resolution of singularities} if the following conditions hold:
	\begin{enumerate}
		\item[(RS1)]	For any $\mathcal{X}\in \ulMCor_k$, there is an abstract blow-up (\Cref{def:abs_blow_up}) $\mathcal{Y}\to \mathcal{X}$ with $\mathcal{Y}\in \ulMCor_k^\ls$.
		\item[(RS2)]	Let $\mathcal{X}\in \ulMCor_k^\ls$ and let $\mathcal{Y}\to \mathcal{X}$ be an abstract blow-up.
						Then there is a sequence of SNC blow-ups $\mathcal{X}_n\to \mathcal{X}_{n-1}\to\dots\to \mathcal{X}_0=\mathcal{X}$ whose composition factors through $\mathcal{Y}$.
	\end{enumerate}
\end{definition}

For example, $k$ admits resolution of singularities when $\ch(k)=0$ \cite{Hironaka1964}.


\begin{corollary}\label{cor:realization}
	Assume that $k$ admits resolution of singularities.
	\begin{enumerate}
		\item	(\cite[Theorem 1.3]{KellyMiyazaki_Hodge2}) For any $q\geq 0$ and $\mathcal{X}\in \ulMCor^\ls_k$, we have an equivalence
				$$
					\Map_{\ulMDM^\eff_k}\bigl(M(\mathcal{X}),\MOmega^q\bigr) \simeq \mathrm{R}\Gamma(X,\MOmega^q_\mathcal{X}).
				$$
				In particular, for any $X\in \Sm_k$ we have
				$$
					\Map_{\ulMDM^\eff_k}\bigl(M(X,\emptyset),\MOmega^q\bigr) \simeq \mathrm{R}\Gamma(X,\Omega^q_X).
				$$
		\item	Suppose that $\ch(k)=p>0$.
				For any $n>0$ and $\mathcal{X}\in \ulMCor^\ls_k$, we have an equivalence
				$$
					\Map_{\ulMDM^\eff_k}\bigl(M(\mathcal{X}),\MW_n\bigr) \simeq \mathrm{R}\Gamma\bigl(X,(\MW_n)_\mathcal{X}\bigr).
				$$
				In particular, for any $X\in \Sm_k$ we have
				$$
					\Map_{\ulMDM^\eff_k}\bigl(M(X,\emptyset),\MW_n\bigr) \simeq \mathrm{R}\Gamma(X,\mathrm{W}_n\mathcal{O}_X).
				$$
	\end{enumerate}
\end{corollary}

\subsection*{Acknowledgments}
The author would like to thank Shane Kelly, Hiroyasu Miyazaki, and Shuji Saito for helpful comments on a draft of this paper.

\subsection*{Conventions}
Throughout this paper, we fix a perfect field $k$.
An \emph{algebraic $k$-scheme} is a separated $k$-scheme of finite type.
For an algebraic $k$-scheme $X$, we write $X^N$ for its normalization.
We write $\Sm_k$ for the category of smooth algebraic $k$-schemes.
We say that a morphism of schemes $f\colon X\to Y$ is \emph{pseudo-dominant} if it sends generic points of $X$ to generic points of $Y$.
If $X$ is a smooth algebraic $k$-scheme, then a \emph{coordinate} of $X$ is a family of functions $x_1,\dots,x_d\in \Gamma(X,\mathcal{O}_X)$ which defines an \'etale morphism $X\to \mathbb{A}^d$.
Unless otherwise specified, cohomologies are taken in the Nisnevich topology.
For a field $K$, we write $K\{t_1,\dots,t_d\}$ for the henselization of $K[t_1,\dots,t_d]$ at $(t_1,\dots,t_d)$.
For a Grothendieck abelian category $\mathcal{A}$, we write $\mathcal{D}(\mathcal{A})$ for its derived $\infty$-category.


\section{Modulus sheaves}

In this section, we recall several definitions and results on modulus sheaves.

\subsection{Finite correspondences}

First we recall Suslin-Voevodsky's category of finite correspondences $\Cor_k$.
It is an additive category having the same objects as $\Sm_k$, and its group of morphisms $\Cor_k(X,Y)$ is the group of algebraic cycles on $X\times Y$ whose components are finite pseudo-dominant over $X$.
Fiber product over $k$ gives a symmetric monoidal structure on $\Cor_k$.
For any morphism $f\colon X\to Y$ in $\Sm_k$, the graph of $f$ gives a morphism $f\colon X\to Y$ in $\Cor_k$.
If $f$ is finite pseudo-dominant, then the transpose of the graph of $f$ gives a morphism ${}^tf\colon Y\to X$ in $\Cor_k$.

We say that a presheaf $F\colon \Cor_k^\op\to \Ab$ is a \emph{Nisnevich sheaf} if for any $X\in \Sm_k$, the presheaf $F_X\colon U\mapsto F(U)$ on $X_\Nis$ is a Nisnevich sheaf.
We write $\Sh_\Nis(\Cor_k)$ for the category of Nisnevich sheaves on $\Cor_k$.
If $f\colon X\to Y$ is a finite pseudo-dominant morphism in $\Sm_k$ and $F\in \Sh_\Nis(\Cor_k)$, then we write $f_*$ or $\Tr_{X/Y}$ for the map $({}^tf)^*\colon F(X)\to F(Y)$.


\subsection{Modulus pairs}

Next we recall basic facts about modulus pairs \cite{KMSY1}.
The definition of modulus pairs is given in the introduction.

An \emph{ambient morphism} of modulus pairs $f\colon \mathcal{X}\to \mathcal{Y}$ is a morphism of $k$-schemes $f\colon X\to Y$ such that $f(X^\circ)\subset Y^\circ$ and $D_X\geq f^*D_Y$ hold.
It is called \emph{minimal} if $D_X=f^*D_Y$ holds.

Let $\mathcal{X},\mathcal{Y}$ be modulus pairs over $k$ and let $V\subset X^\circ\times Y^\circ$ be an integral closed subscheme.
We say that $V$ is \emph{left proper} if the closure $\overline{V}$ of $V$ in $X\times Y$ is proper over $X$.
We say that $V$ is \emph{admissible} if $(\pr_1^*D_X)|_{\overline{V}^N}\geq (\pr_2^*D_Y)|_{\overline{V}^N}$ holds.
We write $\ulMCor_k(\mathcal{X},\mathcal{Y})$ for the subgroup of $\Cor_k(X^\circ,Y^\circ)$ consisting of cycles whose components are left proper and admissible.
This defines a category $\ulMCor_k$ of modulus pairs over $k$.
For any ambient morphism $f\colon \mathcal{X}\to \mathcal{Y}$ over $k$, the graph of $f^\circ$ gives a morphism $f\colon \mathcal{X}\to \mathcal{Y}$ in $\ulMCor_k$.
If $f$ is a proper minimal ambient morphism such that $f^\circ$ is finite pseudo-dominant, then the transpose of the graph of $f^\circ$ gives a morphism ${}^tf\colon \mathcal{Y}\to \mathcal{X}$ in $\ulMCor_k$.
The category $\ulMCor_k$ has a symmetric monoidal structure $\otimes$ given by
$$
	\mathcal{X}\otimes \mathcal{Y}=(X\times Y, \pr_1^*D_X+\pr_2^*D_Y).
$$
\begin{definition}\label{def:abs_blow_up}
	An \emph{abstract blow-up} is a proper minimal ambient morphism $\pi\colon \mathcal{Y}\to \mathcal{X}$ which induces an isomorphism $\pi^\circ\colon Y^\circ\to X^\circ$.
Any such morphism is an isomorphism in $\ulMCor_k$, because the transpose ${}^t\pi\colon \mathcal{X}\to \mathcal{Y}$ gives an inverse.
\end{definition}

Let $\mathcal{X}\in \ulMCor_k^\ls$, i.e.\,$X\in \Sm_k$ and $|D_X|$ has strict normal crossings.
We say that a smooth closed subscheme $Z\subset X$ has \emph{strict normal crossings} to $|D_X|$ if Zariski locally on $X$, there is a coordinate $x_1,\dots,x_d$ such that we can write $|D_X|=\{x_1x_2\cdots x_s=0\}$, $Z=\{x_{i_1}=\cdots=x_{i_r}=0\}$.

\begin{definition}\label{def:SNC_blowup}
	Let $\mathcal{X}\in \ulMCor_k^\ls$.
	An \emph{SNC blow-up} $\mathcal{Y}\to \mathcal{X}$ is an abstract blow-up such that there exists a smooth closed subscheme $Z\subset |D_X|$ which has strict normal crossings to $|D_X|$ and an isomorphism $Y\simeq \Bl_Z(X)$ over $X$.
\end{definition}


We say that a presheaf $F\colon (\ulMCor_k)^\op\to \Ab$ is a \emph{modulus sheaf} if for any $\mathcal{X}\in \ulMCor_k$, the presheaf $F_\mathcal{X}\colon U\mapsto F(U,D_X|_U)$ on $X_\Nis$ is a sheaf.
For example, the representable presheaf $\mathbb{Z}_\tr(\mathcal{X})=\ulMCor_k({-},\mathcal{X})$ is a modulus sheaf \cite[Proposition 3.3.5]{KMSY1}.
We write $\Sh_\Nis(\ulMCor_k)$ for the category of modulus sheaves.
The inclusion functor $\Sh_\Nis(\ulMCor_k)\to \PSh(\ulMCor_k)$ admits an exact left adjoint \cite[Theorem 4.5.5]{KMSY1}.
In particular, $\Sh_\Nis(\ulMCor_k)$ is Grothendieck abelian.




\subsection{$\mathbb{Q}$-divisors}

Following \cite{Koizumi-Miyazaki}, we enlarge the category of modulus pairs by allowing $\mathbb{Q}$-divisors.
A \emph{$\mathbb{Q}$-Cartier divisor} $D$ on a scheme $X$ is an element of $\mathbb{Q}\otimes_\mathbb{Z}\CDiv(X)$ where $\CDiv(X)$ is the group of Cartier divisor of $X$.
It is called \emph{$\mathbb{Q}$-effective} if it is contained in the $\mathbb{Q}_{>0}$-span of some effective Cartier divisor. 
A \emph{$\mathbb{Q}$-modulus pair} over $k$ is a pair $\mathcal{X}=(X,D_X)$ where $X$ is an algebraic $k$-scheme and $D_X$ is a $\mathbb{Q}$-effective $\mathbb{Q}$-Cartier divisor on $X$ such that $X^\circ=X\setminus|D_X|$ is smooth over $k$.
We define a category $\QCor_k$ of $\mathbb{Q}$-modulus pairs in the same way as $\ulMCor_k$ \cite[Definition 2.9]{Koizumi-Miyazaki}.
We can regard $\ulMCor_k$ as a full subcategory of $\QCor_k$.





\subsection{Ramification filtrations}
In this subsection we present a method to construct modulus sheaves from \emph{ramification filtrations}.
This is essentially due to R\"ulling-Saito \cite{RS21}, but we make a small modification to allow $\mathbb{Q}$-divisors.

A \emph{geometric henselian DVF} over $k$ is a discrete valuation field $L$ which is isomorphic to $\mathcal{O}^h_{X,x}$ for some $X\in \Sm_k$ and some point $x\in X$ of codimension $1$.
We write $\Phi_k$ for the class of all geometric henselian DVFs over $k$.
For $L\in \Phi_k$, we write $v_L$ for the normalized valuation, $\mathcal{O}_L$ for the valuation ring, and $\mathfrak{m}_L$ for the maximal ideal of $\mathcal{O}_L$.
Note that if $L\in \Phi_k$ and $K$ is the residue field of $\mathcal{O}_L$, then we have $\mathcal{O}_L\simeq K\{t\}$ for a uniformizer $t\in \mathcal{O}_L$.

\begin{definition}\label{def:ram_fil}
	Let $F\in \Sh_\Nis(\Cor_k)$.
	A \emph{ramification filtration} $\Fil$ on $F$ is collection of increasing filtrations $\{\Fil_r F(L)\}_{r\in \mathbb{Q}_{\geq 0}}$ on $F(L)$ indexed by $L\in \Phi_k$ which satisfies the following:
	\begin{enumerate}
		\item	For any $L\in \Phi_k$, we have $\Im\bigl(F(\mathcal{O}_L)\to F(L)\bigr)\subset \Fil_0F(L)$.
		\item	If $L\in \Phi_k$ and $L'/L$ is a finite extension with ramification index $e$, then we have
				$$
					\Tr_{L'/L}\bigl(\Fil_r F(L')\bigr)\subset \Fil_{\frac{r}{e}}F(L)\quad (r\in \mathbb{Q}_{\geq 0}).
				$$
	\end{enumerate}
\end{definition}

\begin{definition}
	Let $F\in \Sh_\Nis(\Cor_k)$ and let $\Fil$ be a ramification filtration on $F$.
	Let $\mathcal{X}\in \QCor_k$ and $a\in F(X^\circ)$.
	We say that $a$ is \emph{bounded by $D_X$} if for any $L\in \Phi_k$ and any $\rho\colon \Spec L\to X^\circ$ which extends to $\overline{\rho}\colon \Spec \mathcal{O}_L\to X$, we have $\rho^*a\in \Fil_{v_L(\overline{\rho}^*D_X)}F(L)$.
	We write $F_{\Fil}(\mathcal{X})$ for the subgroup of $F(X^\circ)$ consisting of elements bounded by $D_X$.
\end{definition}

\begin{lemma}[cf.\,{\cite[Proposition 4.7]{RS21}}]\label{lem:ram_fil}
	Let $F\in \Sh_\Nis(\Cor_k)$ and let $\Fil$ be a ramification filtration on $F$.
	\begin{enumerate}
		\item	For any $\mathcal{X},\mathcal{Y}\in \QCor_k$ and $\alpha\in \QCor_k(\mathcal{X},\mathcal{Y})$, we have $\alpha^*F_{\Fil} (\mathcal{Y})\subset F_{\Fil} (\mathcal{X})$.
		\item	$F_{\Fil}$ is a modulus sheaf.
	\end{enumerate}
\end{lemma}

\begin{proof}
	(1) We may assume that $\alpha = [V]$ for some integral closed subscheme $V\subset X^\circ\times Y^\circ$.
	Let $b\in F_{\Fil}(\mathcal{Y})$ and $L\in \Phi_k$.
	Suppose that $\rho\colon \Spec L\to X^\circ$ is a morphism which extends to $\overline{\rho}\colon \Spec \mathcal{O}_L\to X$.
	It suffices to show that $\rho^*\alpha^*b\in \Fil_{v_L(\overline{\rho}^*D_X)}F(L)$.
	
	If we set $S=\Spec \mathcal{O}_L$ and $D_S=\overline{\rho}^*D_X$, then the graph of $\rho$ gives a morphism $\mathcal{S}\to \mathcal{X}$.
	We can write $V^N\times_X\Spec L=\coprod_{i=1}^m \Spec L_i$ with $L_1,\dots,L_m\in \Phi_k$.
	We set $S_i=\Spec \mathcal{O}_{L_i}$ and $D_{S_i}=D_S|_{S_i}$.
	Since $\overline{V}^N$ is proper over $X$, the canonical morphism $\Spec L_i\to V$ extends to $S_i \to \overline{V}^N$.
	Moreover, since $V$ is admissible, we get canonical ambient morphisms $\rho_i\colon \mathcal{S}_i\to \mathcal{Y}$ and $\pi_i\colon \mathcal{S}_i\to \mathcal{S}$.
	The composition $\alpha\circ\rho\colon \mathcal{S}\to \mathcal{Y}$ is equal to $\sum_{i=1}^m \rho_i\circ {}^t\pi_i$, so we have $\rho^*\alpha^*b=\sum_{i=1}^m\pi_{i,*}\rho_i^*b$.
	Since $b\in F_{\Fil}(\mathcal{Y})$, we have $\rho_i^*b\in \Fil_{v_{L_i}(D_{S_i})}F(L_i)$.
	By definition of ramification filtration, we get $\pi_{i,*}\rho_i^*b\in \Fil_{v_L(D_S)}F(L)$, as desired.
	
	(2) Since $F_{\Fil}\subset F$, it is clear that $F_{\Fil}$ is separated for the Nisnevich topology.
	Let $\mathcal{X}\in \QCor_k$ and let $\{\pi_i\colon U_i\to X\}_{i\in I}$ be a Nisnevich covering of $X$.
	It suffices to show that if $a\in F(X^\circ)$ satisfies $a|_{U_i}\in F_{\Fil}(U_i,D_X|_{U_i})$ for all $i\in I$, then $a\in F_{\Fil}(\mathcal{X})$.
	Suppose that $\rho\colon \Spec L\to X^\circ$ is a morphism which extends to $\overline{\rho}\colon \Spec \mathcal{O}_L\to X$.
	Since $L$ is henselian, $\overline{\rho}$ lifts to $\overline{\rho}'\colon \Spec \mathcal{O}_L\to U_i$ for some $i$.
	By $a|_{U_i}\in F_{\Fil}(U_i,D_X|_{U_i})$, we get $\rho^*a\in \Fil_{v_L(\overline{\rho}^*D_X)}F(L)$, as desired.
\end{proof}









\section{Examples of modulus sheaves}

In this section, we construct two examples of modulus sheaves.

\subsection{Differential forms}

Fix an integer $q\geq 0$.
Recall from \cite[Corollary A.6.1]{KSY1} that the sheaf of differential forms $\Omega^q=\bigl(X\mapsto \Gamma(X,\Omega^q_X)\bigr)$ can be regarded as an object of $\Sh_\Nis(\Cor_k)$.
Kelly-Miyazaki \cite[Theorem 4.3]{KellyMiyazaki_Hodge2} extended this to a sheaf $\MOmega^q$ on $\ulMCor_k^\ls$ satisfying
$$
	\MOmega^q(\mathcal{X})=\Gamma\bigl(X,\Omega^q_X(\log|D_X|)(D_X-|D_X|)\bigr).
$$
In this subsection, we further extend $\MOmega^q$ to a sheaf on $\QCor_k$ and show that it has cohomological ls cube invariance, affine vanishing property, and left continuity.



\begin{definition}
	Let $L\in \Phi_k$.
	We write $\Omega^q(\mathcal{O}_L)(\logrm)$ for the subgroup of $\Omega^q(L)$ generated by $\Omega^q(\mathcal{O}_L)$ and $\Omega^{q-1}(\mathcal{O}_L)\wedge\dlog(t)$, where $t\in \mathcal{O}_L$ is a uniformizer.
	We define the \emph{logarithmic filtration} on $\Omega^q(L)$ by
	$$
		\Fil_r \Omega^q(L)=\begin{cases}
			\Omega^q(\mathcal{O}_L)&(r=0)\\
			t^{-\ceil{r}+1}\Omega^q(\mathcal{O}_L)(\logrm)&(r>0).
		\end{cases}
	$$
	This does not depend on the choice of the uniformizer $t$.
\end{definition}

\begin{lemma}\label{omega_max}
	Let $L\in \Phi_k$ and let $t\in \mathcal{O}_L$ be a uniformizer.
	For $\omega\in \Omega^q(\mathcal{O}_L)$, we have $\omega\in t\cdot \Omega^q(\mathcal{O}_L)(\logrm)$ if and only if $\dlog(t)\wedge\omega\in \Omega^{q+1}(\mathcal{O}_L)$.
\end{lemma}

\begin{proof}
	The `only if' part is obvious.
	Let $K$ be the residue field of $L$.
	We have an isomorphism $\mathcal{O}_L\simeq K\{t\}$, so we have $\Omega^1(\mathcal{O}_L)\simeq \bigl(\mathcal{O}_L\otimes_K\Omega^1(K)\bigr)\oplus \mathcal{O}_L dt$ and hence
	$$
		\Omega^q(\mathcal{O}_L)\simeq \bigl(\mathcal{O}_L\otimes_K\Omega^q(K)\bigr)\oplus \bigl(\mathcal{O}_L\otimes_K\Omega^{q-1}(K)\bigr)\wedge dt.
	$$
	If we write $\omega = \alpha+\beta\wedge dt\in \Omega^q(\mathcal{O}_L)$, then we have $\dlog(t)\wedge\omega = t^{-1}\alpha\wedge dt$, so $\dlog(t)\wedge\omega\in \Omega^{q+1}(\mathcal{O}_L)$ is equivalent to $\alpha\in t\cdot \Omega^q(K)$, which is equivalent to $\omega\in t\cdot \Omega^q(\mathcal{O}_L)(\logrm)$.
\end{proof}

\begin{lemma}\label{omega_trace}
	Let $L\in \Phi_k$ and let $t\in \mathcal{O}_L$ be a uniformizer.
	Let $L'/L$ be a finite extension and let $t'\in \mathcal{O}_{L'}$ be a uniformizer.
	Then $\Tr_{L'/L}\bigl(t'\cdot \Omega(\mathcal{O}_{L'})(\logrm)\bigr)\subset t\cdot \Omega(\mathcal{O}_L)(\logrm)$.
\end{lemma}

\begin{proof}
	Write $t=u(t')^e$.
	For any $\omega\in \Omega^q(\mathcal{O}_{L'})(\logrm)$, we have
	\begin{align*}
		\dlog(t)\wedge\Tr_{L'/L}(t'\cdot \omega)&{}=\Tr_{L'/L}\bigl(\dlog(u(t')^e)\wedge t'\cdot \omega\bigr)\\
		&{}=\Tr_{L'/L}\bigl(\dlog(u)\wedge t'\cdot \omega\bigr) + e\cdot\Tr_{L'/L}\bigl(\dlog(t')\wedge t'\cdot \omega\bigr)\\
		&{}\in \Tr_{L'/L}\bigl(\Omega^{q+1}(\mathcal{O}_{L'})\bigr)\subset \Omega^{q+1}(\mathcal{O}_L).
	\end{align*}
	By \Cref{omega_max}, this implies that $\Tr_{L'/L}(t'\cdot \omega)\in t\cdot \Omega^q(\mathcal{O}_L)(\logrm)$.
\end{proof}

\begin{lemma}\label{omega_trace_axiom}
	Let $L\in \Phi_k$ and let $L'/L$ be a finite extension of ramification index $e$.
	Then we have $\Tr_{L'/L}\bigl(\Fil_r\Omega^q(L)\bigr)\subset \Fil_{\frac{r}{e}}\Omega^q(L)$ for any $r\in \mathbb{Q}_{\geq 0}$.
\end{lemma}

\begin{proof}
	We may assume that $r>0$.
	Take uniformizers $t\in \mathcal{O}_L$, $t'\in \mathcal{O}_{L'}$ and write $t=u(t')^e$.
	Let $\omega\in \Fil_r \Omega^q(L')$.
	Then we have $t'^{\ceil{r}}\omega\in t'\cdot \Omega^q(\mathcal{O}_{L'})(\logrm)$.
	Using $e\ceil{\frac{r}{e}}\geq r$ and \Cref{omega_trace}, we get
	\begin{align*}
		t^{\ceil{\frac{r}{e}}}\Tr_{L'/L}(\omega)&=\Tr_{L'/L}(u^{\ceil{\frac{r}{e}}}t'^{e\ceil{\frac{r}{e}}}\omega)\\
		&\in \Tr_{L'/L}\bigl(t'\cdot \Omega^q(\mathcal{O}_{L'})(\logrm)\bigr)\subset t\cdot \Omega^q(\mathcal{O}_L)(\logrm).
	\end{align*}
	This implies $\Tr_{L'/L}\bigl(\Fil_r\Omega^q(L)\bigr)\subset \Fil_{\frac{r}{e}}\Omega^q(L)$.
\end{proof}

We write $\MOmega^q$ for the modulus sheaf associated to the logarithmic filtration on $\Omega^q$.
The next lemma shows that this coincides with Kelly-Miyazaki's definition on $\ulMCor_k^\ls$:

\begin{lemma}\label{omega_log_formula}
	For any $\mathcal{X}\in \QCor_k^\ls$, we have
	$$
		\MOmega^q(\mathcal{X})=\Gamma\bigl(X,\Omega^q_X(\log |D_X|)(\ceil{D_X}-|D_X|)\bigr).
	$$
	In particular, $\MOmega^q$ has left continuity and affine vanishing property.
\end{lemma}

\begin{proof}
	It is clear that the left hand side is contained in the right hand side.
	Let us prove the inverse inclusion.
	We may assume that $\mathcal{X}$ has a coordinate $x_1,\dots,x_d$ such that $D_X=\sum_{i=1}^s r_i D_i$ where $D_i=\{x_i=0\}$.
	Let $\omega$ be an element of the right hand side of the stated formula.
	Let $L\in \Phi_k$ and let $t\in \mathcal{O}_L$ be a uniformizer.
	Suppose that $\rho\colon \Spec L\to X^\circ$ is a morphism which extends to $\overline{\rho}\colon \Spec \mathcal{O}_L\to X$.
	We set $e_i=v_L(\rho^*x_i)$ and suppose that $e_i>0$ for some $i$.
	Then there is some unit $u\in \mathcal{O}_L^\times$ such that
	$$
		u t^{\sum_{i=1}^s e_i(\ceil{r_i}-1)}\rho^*\omega = \rho^*(x_1^{\ceil{r_1}-1}\cdots x_s^{\ceil{r_s}-1}\omega)\in \Omega^q(\mathcal{O}_L)(\logrm).
	$$
	On the other hand, we have $\sum_{i=1}^s e_i(\ceil{r_i}-1)<\sum_{i=1}^s e_i r_i\leq \ceil{\sum_{i=1}^s e_i r_i}$, so we get
	$$
		t^{\ceil{\sum_{i=1}^s e_i r_i}-1}\rho^*\omega\in \Omega^q(\mathcal{O}_L)(\logrm).
	$$
	This shows that $\omega\in \MOmega^q(\mathcal{X})$.
\end{proof}

\begin{lemma}\label{omega_CCI}
	The modulus sheaf $\MOmega^q$ has cohomological ls cube invariance.
\end{lemma}

\begin{proof}
	This follows from \Cref{omega_log_formula} and \cite[Corollary 5.2]{KellyMiyazaki_Hodge2}.
	We sketch the proof for readers' convenience.
	Let $\mathcal{X}\in \QCor_k^\ls$.
	An explicit computation shows that 
	\begin{align*}
		(\MOmega^q)_{\mathcal{X}\otimes\bcube}&{}\simeq \pr_1^*(\MOmega^q)_\mathcal{X}\oplus\bigl(\pr_1^*(\MOmega^{q-1})_\mathcal{X}\otimes\pr_2^*(\MOmega^1)_\bcube\bigr)\\
		&{}\simeq \pr_1^*(\MOmega^q)_\mathcal{X}\oplus\bigl(\pr_1^*(\MOmega^{q-1})_\mathcal{X}\otimes\pr_2^*\mathcal{O}_{\mathbb{P}^1}(-1)\bigr)
	\end{align*}
	as a quasi-coherent sheaf on $X\times\mathbb{P}^1$.
	Therefore it suffices to show that for any finite locally free $\mathcal{O}_X$-modules $E,E'$ on $X$, we have $\mathrm{R}\pr_{1,*}\bigl(\pr_1^*E\oplus(\pr_1^*E'\otimes \pr_2^*\mathcal{O}_{\mathbb{P}^1}(-1))\bigr)\simeq E$.
	This can be reduced to the case $E=E'=\mathcal{O}_X$, for which the assertion is obvious.
\end{proof}


\subsection{Witt vectors}
In this subsection we assume that $\ch(k)=p>0$.
Fix an integer $n>0$.
The sheaf of $p$-typical Witt vectors $\mathrm{W}_n=\bigl(X\mapsto \Gamma(X,\mathrm{W}_n\mathcal{O}_X)\bigr)$ is representable by an algebraic group, so it can be regarded as an object of $\Sh_\Nis(\Cor_k)$ \cite[Lemma 1.4.4]{BVK16}.
In this subsection, we construct a modulus sheaf $\MW_n\in \Sh_\Nis(\QCor_k)$ extending $\mathrm{W}_n$ which has cohomological ls cube invariance, affine vanishing property, and left continuity.

\begin{definition}
	Let $L\in \Phi_k$ and let $t\in \mathcal{O}_L$ be a uniformizer.
	We define the \emph{Brylinski-Kato filtration} \cite[Definition 3.1]{Kato89} on $\mathrm{W}_n(L)$ by
	$$
		\Fil_r \mathrm{W}_n(L)=\begin{cases}
			\mathrm{W}_n(\mathcal{O}_L)&(r=0)\\
			\bigl\{a\in \mathrm{W}_n(L)\bigm| [t^{\ceil{r}-1}]\mathrm{F}^{n-1}(a)\in \mathrm{W}_n(\mathcal{O}_L)\bigr\}&(r>0).
		\end{cases}
	$$
	This does not depend on the choice of the uniformizer $t$.
\end{definition}

\begin{remark}\label{Witt_filtration_valuation}
	For $a=(a_0,\dots,a_{n-1})\in \mathrm{W}_n(L)$, we have
	$$
		[t^{\ceil{r}-1}]\mathrm{F}^{n-1}(a)=(t^{\ceil{r}-1}a_0^{p^{n-1}},t^{p(\ceil{r}-1)}a_1^{p^{n-1}},\dots,t^{p^{n-1}(\ceil{r}-1)}a_{n-1}^{p^{n-1}}).
	$$
	Therefore we have $a\in \Fil_r \mathrm{W}_n(L)$ if and only if
	$$
		v_L(a_i) +  p^i\cdot \dfrac{\ceil{r}-1}{p^{n-1}} \geq 0
	$$
	holds for $i=0,1,\dots,n-1$.
	In particular, we have $\Fil_1\mathrm{W}_n(L)=\Fil_0\mathrm{W}_n(L)$.
\end{remark}

\begin{lemma}\label{Witt_trace}
	Let $L\in \Phi_k$.
	Let $L'/L$ be a finite extension and $t'\in \mathcal{O}_{L'}$ be a uniformizer.
	If $\overline{L}$ is an algebraic closure of $L$, then we have
	$$i_{\overline{L}/L}\circ \Tr_{L'/L}\bigl([t']\cdot \mathrm{W}_n(\mathcal{O}_{L'})\bigr)\subset [t']\cdot \mathrm{W}_n(\mathcal{O}_{\overline{L}})$$
	where $i_{\overline{L}/L}\colon \mathrm{W}_n(\mathcal{O}_L)\to \mathrm{W}_n(\mathcal{O}_{\overline{L}})$ is the inclusion.
\end{lemma}

\begin{proof}
	Let $a\in \mathrm{W}_n(\mathcal{O}_{L'})$.
	If $m$ (resp. $n$) is the separable (resp. inseparable) degree of $L'/L$, then there are $m$ distinct $L$-embeddings $\sigma_1,\dots,\sigma_m\colon L'\to \overline{L}$ and we have
	$$
		i_{\overline{L}/L}\circ \Tr_{L'/L}([t']\cdot a)=n\cdot \sum_{i=1}^m\sigma_i([t']\cdot a).
	$$
	It is clear that the right hand side is contained in $[t']\cdot \mathrm{W}_n(\mathcal{O}_{\overline{L}})$.
\end{proof}

\begin{lemma}\label{Witt_trace_axiom}
	Let $L\in \Phi_k$ and let $L'/L$ be a finite extension of ramification index $e$.
	Then we have $\Tr_{L'/L}(\Fil_r \mathrm{W}_n(L))\subset \Fil_{\frac{r}{e}}\mathrm{W}_n(L)$ for any $r\in \mathbb{Q}_{\geq 0}$.
\end{lemma}

\begin{proof}
	Take uniformizers $t\in \mathcal{O}_L$, $t'\in \mathcal{O}_{L'}$ and write $t=u(t')^e$.
	Let $a\in \Fil_r \mathrm{W}_n(L')$.
	Then we have $[t'^{\ceil{r}-1}]\mathrm{F}^{n-1}(a)\in \mathrm{W}_n(\mathcal{O}_{L'})$.
	Using $e\ceil{\frac{r}{e}}\geq r$ and \Cref{Witt_trace}, we get
	\begin{align*}
		i_{\overline{L}/L}\bigl([t^{\ceil{\frac{r}{e}}}]\mathrm{F}^{n-1}(\Tr_{L'/L}(a))\bigr)
		&=i_{\overline{L}/L}\circ\Tr_{L'/L}\bigl([u^{\ceil{\frac{r}{e}}}t'^{e\ceil{\frac{r}{e}}}]\mathrm{F}^{n-1}(a)\bigr)\\
		&\in i_{\overline{L}/L}\circ \Tr_{L'/L}\bigl([t']\cdot \mathrm{W}_n(\mathcal{O}_{L'})\bigl)
		\subset [t']\cdot \mathrm{W}_n(\mathcal{O}_{\overline{L}}).
	\end{align*}
	If we set $\Tr_{L'/L}(a)=b=(b_0,\dots,b_{n-1})$, then the above computation shows that
	$$
		\biggl\lceil\frac{r}{e}\biggr\rceil+p^{n-1-i}v_L(b_i)\geq\dfrac{1}{e}
	$$
	for $i=0,1,\dots,n-1$.
	Therefore the left hand side is a positive integer, so we get
	$$
		v_L(b_i)+p^i\cdot\frac{\ceil{\frac{r}{e}}-1}{p^{n-1}}\geq 0
	$$
	for $i=0,1,\dots,n-1$.
	This implies $\Tr_{L'/L}(a)=b\in \Fil_{\frac{r}{e}}\mathrm{W}_n(L)$.
\end{proof}

We write $\MW_n$ for the modulus sheaf associated to the Brylinski-Kato filtration on $\mathrm{W}_n$.
In order to give an explicit formula for $\MW_n$, we recall the notion of \emph{Witt divisorial sheaf} which was introduced by Tanaka \cite{Tanaka_Witt}.

\begin{definition}[{\cite[Definition 3.1]{Tanaka_Witt}}]
	Let $X\in \Sm_k$ and let $D=\sum_{i=1}^s r_i D_i$ be a $\mathbb{Q}$-divisor on $X$.
	Let $\xi_i$ denote the generic point of $D_i$ and $v_i$ denote the discrete valuation associated to $D_i$.
	Let $j\colon X\setminus|D|\to X$ be the inclusion.
	The \emph{Witt divisorial sheaf} $\mathrm{W}_n\mathcal{O}_X(D)$ is the Zariski subsheaf of $j_*\mathrm{W}_n\mathcal{O}_{X\setminus |D|}$ defined by
	$$
		U\mapsto \bigl\{(a_0,\dots,a_{n-1})\in \Gamma(U, j_*\mathrm{W}_n\mathcal{O}_{X\setminus|D|})\bigm| \xi_i\in U\implies v_i(a_i)+p^i r_i\geq 0\bigr\}.
	$$
	It is known that $\mathrm{W}_n\mathcal{O}_X(D)$ is a quasi-coherent $\mathrm{W}_n\mathcal{O}_X$-module \cite[Lemma 3.5 (2)]{Tanaka_Witt}.
\end{definition}

\begin{lemma}\label{Witt_log_formula}
	For any $\mathcal{X}\in \QCor_k^\ls$, we have
	$$
		\MW_n(\mathcal{X})=\Gamma\biggl(X,\mathrm{W}_n\mathcal{O}_X\biggl(\frac{\ceil{D_X}-|D_X|}{p^{n-1}}\biggr)\biggr).
	$$
	In particular, $\MW_n$ has left continuity and affine vanishing property.
\end{lemma}

\begin{proof}
	It is clear from \Cref{Witt_filtration_valuation} that the left hand side is contained in the right hand side.
	Let us prove the inverse inclusion.
	We may assume that $\mathcal{X}$ has a coordinate $x_1,\dots,x_d$ such that $D_X=\sum_{i=1}^s r_i D_i$ where $D_i=\{x_i=0\}$.
	Let $a$ be an element of the right hand side of the stated formula.
	Let $L\in \Phi_k$ and let $t\in \mathcal{O}_L$ be a uniformizer.
	Suppose that $\rho\colon \Spec L\to X^\circ$ is a morphism which extends to $\overline{\rho}\colon \Spec \mathcal{O}_L\to X$.
	We set $e_i=v_L(\rho^*x_i)$ and suppose that $e_i>0$ for some $i$.
	Then there is some unit $u\in \mathcal{O}_L^\times$ such that
	\begin{align*}
		[u t^{\sum_{i=1}^s e_i(\ceil{r_i}-1)}]\mathrm{F}^{n-1}(\rho^*a)&{} = \rho^*\bigl([x_1^{\ceil{r_1}-1}\cdots x_s^{\ceil{r_s}-1}]\mathrm{F}^{n-1}(a)\bigr)\\
		&{} \in \rho^*\bigl(\Gamma(X,\mathrm{W}_n\mathcal{O}_X)\bigr)\subset \mathrm{W}_n(\mathcal{O}_L).
	\end{align*}
	On the other hand, we have $\sum_{i=1}^s e_i(\ceil{r_i}-1)<\sum_{i=1}^s e_i r_i\leq \ceil{\sum_{i=1}^s e_i r_i}$, so we get
	$$
		[t^{\ceil{\sum_{i=1}^s e_i r_i}-1}]\mathrm{F}^{n-1}(\rho^*a) \in \mathrm{W}_n(\mathcal{O}_L).
	$$
	This shows that $a\in \MW_n(\mathcal{X})$.
\end{proof}

\begin{lemma}\label{Witt_CCI}
	The modulus sheaf $\MW_n$ has cohomological ls cube invariance.
\end{lemma}

\begin{proof}
	Let $X\in \Sm_k$.
	By \Cref{Witt_log_formula}, it suffices to show that the canonical morphism
	$$
		\mathrm{W}_n\mathcal{O}_X(D)\to \mathrm{R}\pr_{1,*}\mathrm{W}_n\mathcal{O}_{X\times\mathbb{P}^1}(\pi^*D)
	$$
	is an equivalence for any $\mathbb{Q}$-divisor $D$ on $X$.
	We prove this by induction on $n$.
	If $n=1$, then what we have to show is that the canonical morphism
	$\mathcal{O}_X(D)\to \mathrm{R}\pr_{1,*}\pr_1^*\mathcal{O}_X(D)$
	is an isomorphism, which is obvious.
	Suppose that $n\geq 2$.
	By \Cref{Witt_log_formula} and \cite[Proposition 3.7]{Tanaka_Witt}, we have an exact sequence
	$$
		0\to (\mathrm{F}^{n-1}_X)_*\mathcal{O}_X(p^{n-1}D)\to \mathrm{W}_n\mathcal{O}_X(D)\to \mathrm{W}_{n-1}\mathcal{O}_X(D)\to 0
	$$
	of Zariski sheaves on $X$, where $\mathrm{F}_X\colon X\to X$ is the absolute Frobenius morphism.
	We also have a similar exact sequence on $X\times\mathbb{P}^1$.
	Therefore the claim follows from the induction hypothesis.
\end{proof}






\section{Cohomological blow-up invariance}\label{section_main}



In this section we prove \Cref{main}.

\begin{lemma}\label{P1_vanishing_lemma}
	Let $A$ be a henselian local ring.
	Let $\varphi\colon F\to G$ be a morphism of Nisnevich sheaves on $\mathbb{P}^1_A$ such that $\Ker \varphi$ and $\Coker \varphi$ are supported on $\{0,\infty\}\subset\mathbb{P}^1_A$.
	Then the morphism
	$$
		\mathrm{H}^i(\mathbb{P}^1_A, F)\to \mathrm{H}^i(\mathbb{P}^1_A, G)
	$$
	is surjective for $i>0$.
\end{lemma}

\begin{proof}
	We have two exact sequences of Nisnevich sheaves
	\begin{align*}
		&0\to \Ker f\to F\xrightarrow{\varphi} \Im f\to 0,\\
		&0\to \Im f\to G\to \Coker f\to 0.
	\end{align*}
	Taking the cohomology, we get exact sequences
	\begin{align*}
		&H^i(\mathbb{P}^1_A,F)\xrightarrow{\varphi} H^i(\mathbb{P}^1_A,\Im f)\to H^{i+1}(\mathbb{P}^1_A,\Ker f),\\
		&H^i(\mathbb{P}^1_A,\Im f)\to H^i(\mathbb{P}^1_A,G)\to H^i(\mathbb{P}^1_A,\Coker f).
	\end{align*}
	By our assumption, the rightmost terms are $0$ for $i>0$, so the composition
	$$
		H^i(\mathbb{P}^1_A,F)\xrightarrow{\varphi} H^i(\mathbb{P}^1_A,\Im f)\to H^i(\mathbb{P}^1_A,G)
	$$
	is surjective for $i>0$.
\end{proof}


The blow-up of $\mathbb{A}^2$ at the origin is isomorphic to the total space of the line bundle $\mathcal{O}(-1)$ on $\mathbb{P}^1$.
We first study the cohomology of this scheme.
Let $H^{(n)}$ denote the total space of the line bundle $\mathcal{O}(n)$ on $\mathbb{P}^1$;
it is obtained by gluing $\Spec k[s,x_n]$ and $\Spec k[t,y_n]$ via
$$
	\Spec k[s^\pm,x_n]\xrightarrow{\sim}\Spec k[t^\pm,y_n];\quad t\mapsto s^{-1},\;y_n\mapsto s^{-n}x_n.
$$
Let $\pi^{(n)}\colon H^{(n)}\to \mathbb{P}^1$ be the canonical projection.
We define divisors $D_0^{(n)}, D_\infty^{(n)}, E^{(n)}$ on $H^{(n)}$ by
$$
	D_0^{(n)} = \{s=0\},\quad D_\infty^{(n)} = \{t=0\}, \quad E^{(n)} = \{x_n=0\}\cup \{y_n=0\}.
$$
We define $\mathcal{H}^{(n)}(a,b,c)=(H^{(n)}, a D_0^{(n)}+b D_\infty^{(n)}+c E^{(n)})$ for $a,b,c\in \mathbb{Q}_{\geq 0}$.
For $\mathcal{X}\in \QCor_k^\ls$, we write $\mathcal{H}^{(n)}_\mathcal{X}(a,b,c) = \mathcal{H}^{(n)}(a,b,c)\otimes \mathcal{X}$ and write $\pi^{(n)}_X\colon H^{(n)}_X\to \mathbb{P}^1_X$ for the base change of $\pi^{(n)}$.

\begin{lemma}\label{Hirzebruch_vanishing}
	Suppose that $F\in \Sh_\Nis(\QCor_k)$ has cohomological ls cube invariance, affine vanishing property, and left continuity.
	Let $\mathcal{X}$ be the henselian localization of an object of $\QCor_k^\ls$ and $a,b\in \mathbb{Q}_{\geq 0}$, $a\neq 0$.
	Then we have
    $$
		\mathrm{H}^i\bigl(H^{(-1)}_X, F_{\mathcal{H}^{(-1)}_\mathcal{X}(a,b,a+b)}\bigr)=0\quad (i>0).
	$$
\end{lemma}

\begin{proof}
	It suffices to prove that for each prime number $\ell$, the cohomology groups are $\ell^\infty$-torsion.
	Since $F$ has left continuity, it suffices to show that for $N=\ell^e$, the groups
	$$
		\mathrm{H}^i\bigl(H^{(-1)}_X, F_{\mathcal{H}^{(-1)}_\mathcal{X}(a,b,a+b-\frac{a+b+1}{N+1})}\bigr)\quad (i>0)
	$$
	are annihilated by $N$.
	Set $c=a+b-\frac{a+b+1}{N+1}$.
	Then $(Na-1)+Nb=(N+1)c$, so there are integers $m,m'>0$ such that
	$$
		m c \leq N a-1,\quad m' c\leq N b,\quad m+m'=N.
	$$
	Let $\rho_{N}\colon \mathbb{P}^1_X\to \mathbb{P}^1_X$ be the morphism $s\mapsto s^N$.
	We have the following Cartesian diagram:
	$$
	\xymatrix{
		H^{(-N)}_X\ar[r]^-{\theta_N}\ar[d]		&H^{(-1)}_X\ar[d]\\
		\mathbb{P}^1_X\ar[r]^-{\rho_N}				&\mathbb{P}^1_X
	}
	$$
	The morphism $\theta_N$ is finite locally free of degree $N$, and it induces a minimal ambient morphism $\mathcal{H}_\mathcal{X}^{(-N)}(Na,Nb,c)\to \mathcal{H}_\mathcal{X}^{(-1)}(a,b,c)$.
	The composition
	$$
		F_{\mathcal{H}^{(-1)}_\mathcal{X}(a,b,c)}\xrightarrow{\theta_N^*} \theta_{N,*} F_{\mathcal{H}^{(-N)}_\mathcal{X}(Na,Nb,c)}\xrightarrow{\theta_{N,*}} F_{\mathcal{H}^{(-1)}_\mathcal{X}(a,b,c)}
	$$
	is equal to the multiplication by $N$, so it suffices to show that the group 
	$$
		\mathrm{H}^i\bigl(H^{(-1)}_X, \theta_{N,*}F_{\mathcal{H}^{(-N)}_\mathcal{X}(Na,Nb,c)}\bigr)\simeq
		\mathrm{H}^i\bigl(H^{(-N)}_X, F_{\mathcal{H}^{(-N)}_\mathcal{X}(Na,Nb,c)}\bigr)
	$$
	is zero for $i>0$.
	We have $\mathrm{R}^i\pi^{(-N)}_{X,*}F_{\mathcal{H}^{(-N)}_\mathcal{X}(Na,Nb,c)}=0$ for $i>0$ by affine vanishing property, so it suffices to prove that $\mathrm{H}^i\bigl(\mathbb{P}^1_X, \pi^{(-N)}_{X,*}F_{\mathcal{H}_\mathcal{X}^{(-N)}(Na,Nb,c)}\bigr)=0$ for $i>0$.
	We define a morphism $\psi_{N,m}\colon H^{(-N)}_X\to H^{(0)}_X$ over $\mathbb{P}^1_X$ by
	$$
		x_0\mapsto s^m x_{-N},\quad y_0\mapsto t^{m'}y_{-N}.
	$$
	Then $\psi_{N,m}$ induces an ambient morphism $\mathcal{H}_\mathcal{X}^{(-N)}(Na,Nb,c)\to \mathcal{H}^{(0)}_\mathcal{X}(1,0,c)$.
	The induced morphism
	$$
		\psi_{N,m}^*\colon \pi^{(0)}_{X,*}F_{\mathcal{H}_\mathcal{X}^{(0)}(1,0,c)} \to \pi^{(-N)}_{X,*}F_{\mathcal{H}_\mathcal{X}^{(-N)}(Na,Nb,c)}
	$$
	is an isomorphism on $\mathbb{P}^1_X\setminus\{0,\infty\}$, so the kernel and the cokernel of $\psi_{N,m}^*$ are supported on $\{0,\infty\}\subset\mathbb{P}^1_X$.
	By \Cref{P1_vanishing_lemma}, this shows that the induced map
	$$
		\mathrm{H}^i\bigl(\mathbb{P}^1_X, \pi^{(0)}_{X,*}F_{\mathcal{H}_\mathcal{X}^{(0)}(1,0,c)}\bigr)\to
		\mathrm{H}^i\bigl(\mathbb{P}^1_X, \pi^{(-N)}_{X,*}F_{\mathcal{H}_\mathcal{X}^{(-N)}(Na,Nb,c)}\bigr)
	$$
	is surjective for $i>0$.
	The source of this map can be computed as
	\begin{align*}
		&\mathrm{H}^i(\mathbb{P}^1\times X, \pr_{13,*}F_{\bcube \otimes (\mathbb{A}^1,c[0])\otimes\mathcal{X}})\\
		\simeq{} &\mathrm{H}^i(\mathbb{P}^1\times\mathbb{A}^1\times X, F_{\bcube \otimes (\mathbb{A}^1,c[0])\otimes\mathcal{X}})
		&\cdots\text{(affine vanishing property)}\\
		\simeq{} &\mathrm{H}^i(\mathbb{A}^1\times X, F_{(\mathbb{A}^1,c[0])\otimes\mathcal{X}})
		&\cdots\text{(cohomological cube invariance)}\\
		={} &0
		&\cdots\text{(affine vanishing property)}
	\end{align*}
	for $i>0$.
	Therefore the target is also zero for $i>0$.
\end{proof}

\begin{lemma}\label{blowup-a2}
	Suppose that $F\in \Sh_\Nis(\QCor_k)$ has cohomological ls cube invariance, affine vanishing property, and left continuity.
	Let $\varphi\colon B\to \mathbb{A}^2$ be the blow-up of $\mathbb{A}^2$ at the origin and $L=\mathbb{A}^1\times\{0\}$, $L'=\{0\}\times\mathbb{A}^1$.
	Fix $a,b\in \mathbb{Q}_{\geq 0}$ with $a\neq 0$ and set $\mathcal{B}=\bigl(B,\varphi^*(a L+b L')\bigr)$.
	Then we have
	$$
		\mathrm{R}^i\varphi_{X,*}F_{\mathcal{B}_\mathcal{X}}=0\quad (i>0)
	$$
	for any $\mathcal{X}\in \QCor_k^\ls$, where $\mathcal{B}_\mathcal{X} =\mathcal{B}\otimes \mathcal{X}$ and $\varphi_X\colon B_X\to \mathbb{A}^2_X$ is the base change of $\varphi$.
\end{lemma}

\begin{proof}
	We may replace $\mathcal{X}$ by its henselian localization.
	The higher direct images $\mathrm{R}^i\varphi_{X,*}F_{\mathcal{B}_\mathcal{X}}\;(i>0)$ are supported on $\{(0,0)\}\subset\mathbb{A}^2_X$ and hence
	$$
		\mathrm{H}^j(\mathbb{A}^2_X, \mathrm{R}^i\varphi_{X,*}F_{\mathcal{B}_\mathcal{X}})=0\quad (i,j>0).
	$$
	On the other hand, we have
	$$
		\mathrm{H}^j(\mathbb{A}^2_X, \varphi_{X,*}F_{\mathcal{B}_\mathcal{X}})=\mathrm{H}^j(\mathbb{A}^2\times X,F_{(\mathbb{A}^2,a L+b L')\otimes \mathcal{X}})=0 \quad (j>0)
	$$
	by affine vanishing property, so the Leray spectral sequence yields an isomorphism
	$$
		\mathrm{H}^0(\mathbb{A}^2_X, \mathrm{R}^i\varphi_{X,*}F_{\mathcal{B}_\mathcal{X}})\simeq \mathrm{H}^i(B_X,F_{\mathcal{B}_\mathcal{X}}).
	$$
	It suffices to show that the right hand side vanishes for $i>0$.
	Now the claim follows from \Cref{Hirzebruch_vanishing} and the fact that $\mathcal{B}_\mathcal{X}$ is isomorphic to $\mathcal{H}_\mathcal{X}^{(-1)}(a,b,a+b)$.
\end{proof}

\begin{lemma}\label{blowup_exchange}
	Let $X$ be an algebraic $k$-scheme and $Z\subset Z'\subset X$ be closed subschemes.
	Set $Y=\Bl_Z(X)$ and $Y'=\Bl_{\widetilde{Z'}}(Y)$ where $\widetilde{Z'}$ is the strict transform of $Z'$.
	Set $W=\Bl_{Z'}(X)$ and $W'=\Bl_{\psi^{-1}(Z)}(W)$.
	Then there is an isomorphism $Y'\simeq W'$ over $X$.
\end{lemma}

\begin{proof}
	See \cite[Lemma 2.15]{BRS}.
\end{proof}






\begin{proof}[\bf{Proof of \Cref{main}}]
	Let $\mathcal{X}\in \QCor_k^\ls$ and let $\varphi\colon \mathcal{Y}\to \mathcal{X}$ be an SNC blow-up with center $Z$.
	We prove $\mathrm{R}\varphi_*F_\mathcal{Y}\simeq F_\mathcal{X}$ by induction on $r=\codim_X(Z)$.
	Since the claim is Nisnevich local, we may assume that $X=\Spec K\{x_1,\dots,x_d\}$, $|D_X|=\{x_1x_2\cdots x_s=0\}$, and $Z=\{x_{i_1}=\cdots=x_{i_r}=0\}$.
	We may assume that $i_1=1$.
	The case $r=1$ is trivial.
	If $r=2$, then the claim follows from \Cref{blowup-a2}.
	If $r\geq 3$, then we set $Z'=\{x_1=x_2=0\}\supset Z$.
	Let $\varphi'\colon \mathcal{Y}'\to \mathcal{Y}$ be the SNC blow-up along the strict transform $\widetilde{Z'}$ of $Z'$, $\psi\colon \mathcal{W}\to \mathcal{X}$ be the SNC blow-up along $Z'$, and let $\psi'\colon \mathcal{W}'\to \mathcal{W}$ be the SNC blow-up along $\psi^{-1}(Z)$.
	Then there is an isomorphism $Y'\simeq W'$ over $X$ by \Cref{blowup_exchange}.
	Since $\codim_Y(\widetilde{Z'})=\codim_X(Z')=2$ and $\codim_W(\psi^{-1}(Z))=r-1$, we have
	$$
		\mathrm{R}\varphi'_*F_\mathcal{Y'}\simeq F_\mathcal{Y},\quad \mathrm{R}\psi_*F_\mathcal{W}\simeq F_\mathcal{X},\quad \mathrm{R}\psi'_*F_{\mathcal{W}'}\simeq F_\mathcal{W}.
	$$
	Combining these results, we obtain $\mathrm{R}\varphi_*F_\mathcal{Y}\simeq F_\mathcal{X}$.
\end{proof}



\section{Realizations of motives with modulus}

We show that the Hodge cohomology and the Witt vector cohomology is representable in the category of motives with modulus under resolution of singularities.
First we recall the construction of the category of motives with modulus.

\begin{definition}
	Let $(\mathrm{CI})$ be the localizing subcategory of $\mathcal{D}\bigl(\Sh_\Nis(\ulMCor_k)\bigr)$ generated by the cofiber of $\pr_1\colon \mathbb{Z}_\tr(\mathcal{X}\otimes \bcube)\to \mathbb{Z}_\tr(\mathcal{X})$ for $\mathcal{X}\in \ulMCor_k$.
	The category of motives with modulus $\ulMDM^\eff$ is defined as the Verdier quotient
	$$
		\ulMDM^\eff_k = \frac{\mathcal{D}\bigl(\Sh_\Nis(\ulMCor_k)\bigr)}{(\mathrm{CI})}
	$$
	in the sense of \cite[Definition 5.4]{BGT}.
	By \cite[Theorem 4.1.1]{KMSY3}, this is equivalent to the original definition \cite[Definition 3.2.4]{KMSY3} after taking the homotopy category.
\end{definition}



\begin{lemma}\label{general_realization}
	Assume that $k$ admits resolution of singularities.
	Suppose that $F\in \Sh_\Nis(\ulMCor_k)$ has cohomological ls cube invariance and cohomological blow-up invariance.
	Then for $\mathcal{X}\in \ulMCor_k^\ls$, we have an equivalence
	$$
		\Map_{\ulMDM^\eff_k}\bigl(M(\mathcal{X}),F\bigr) \simeq \mathrm{R}\Gamma(X,F_\mathcal{X}).
	$$
\end{lemma}

\begin{proof}
	Let $\mathcal{X}\in \ulMCor_k^\ls$.
	By \cite[Theorem 4.6.3]{KMSY1} we have
	\begin{align*}
		\Ext^i_{\Sh_\Nis(\ulMCor_k)}\bigl(\mathbb{Z}_\tr(\mathcal{X}),F\bigr)\simeq \colim_{\mathcal{Y}\to \mathcal{X}} \mathrm{H}^i(Y,F_\mathcal{Y})
	\end{align*}
	where the colimit is taken over abstract blow-ups.
	Since $F$ has cohomological blow-up invariance and $k$ satisfies (RS2), the last term is isomorphic to $\mathrm{H}^i(X,F_\mathcal{X})$.
	Therefore the canonical morphism
	$$
		\mathrm{R}\Gamma(X,F_\mathcal{X})\to \Map_{\mathcal{D}(\Sh_\Nis(\ulMCor_k))} \bigl(\mathbb{Z}_\tr(\mathcal{X}),F\bigr)
	$$
	is an equivalence.
	It remains to show that $F$ is $\bcube$-local, i.e., the morphism
	$$
		\Map_{\mathcal{D}(\Sh_\Nis(\ulMCor_k))} \bigl(\mathbb{Z}_\tr(\mathcal{Y}),F\bigr)\to
		\Map_{\mathcal{D}(\Sh_\Nis(\ulMCor_k))} \bigl(\mathbb{Z}_\tr(\mathcal{Y}\otimes\bcube),F\bigr)
	$$
	is an equivalence for any $\mathcal{Y}\in \ulMCor_k$.
	By (RS1), we may assume that $\mathcal{Y}\in \ulMCor_k^\ls$.
	Then this morphism can be identified with $\mathrm{R}\Gamma(Y,F_\mathcal{Y})\to \mathrm{R}\Gamma(Y\times\mathbb{P}^1,F_\mathcal{Y\otimes\bcube})$, so the claim follows from cohomological ls cube invariance.
\end{proof}

Now \Cref{cor:realization} follows from \Cref{omega_witt_CBI}, \Cref{omega_CCI}, \Cref{Witt_CCI}, and \Cref{general_realization}.



\printbibliography

\end{document}

